% Pursuing Stacks (cotes 134-2 to 134-8): archive cote 134-6, site page 113 (PDF page 114), typescript page 431.
% typescript. Draft compared: ps5.tex:7389-7469.
% AI-transcribed by Claude Opus 5.5 (claude-opus-5-5, adaptive thinking, effort high), Message Batches API, batch ps3 (msgbatch_01XdFmuVVBAcgc7EFxdtCqsj), 2026-11-12; reviewed against both supplied scan halves; see CHECK_DECISIONS.tsv.
\scanpage{134-6/113}
\mspage{431}

and not the pairing $R\underline{Hom}_k$ in (111), and hence no formula (ii) (p.428)
relating the two - which makes me feel a little silly ! I'll have to
come back upon this later. \xout{For the}\tikz[baseline=-.4ex]{\draw[color=blue,line width=.55pt] (0,0)--(.45,0)--(.45,.28)--(.95,.28);\draw[color=blue,line width=.4pt] (.42,0)--(.42,.28);} At present, let's dualize what we
got, assuming now \over[typed]{w}{t}hat \hand{$\mathcal{N}$} is a k-additive abelian category stable under
small \hand{\uline{\textcolor{black}{direct}}} limits, and such that \hand{\uline{\textcolor{black}{a small direct sum of monomorphisms}}}
\hand{\uline{\textcolor{black}{in \hand{$\mathcal{N}$} is again a monomorphism}}}. We thenget a pairing
\[ \begin{aligned} (117)\qquad & (F\hand{.,}L\hand{.'})\,\deleted{\xout{\illegible{1}}}\,\hand{\mapsto}\, F.\overset{\over{L}{L}}{*}_k L.\hand{'} \\
& :\quad \deleted{\xout{D^{\hand{\unsure{A}}}\illegible{1}}}\; D^-(A^\wedge_{\hand{\mathcal{N}}}) \times D^-(B^\wedge_k) \hand{\to} D^-(\hand{\mathcal{N}}) \quad , \end{aligned} \]
defined by the formula (dual to (112))
\[ (118)\qquad F.\overset{L}{*}_k L.' \hand{\simeq} F.*_k \hand{\mathcal{L}.'} \quad , \]
where in the second member \hand{$\mathcal{L}.'$} is a projective resolution of L. in
$B^\wedge_k$ ( $B = A^o$ being of course the dual category of A), and the $*_k$ deno-%
tes the simple complex associated to the double complex obtained by
applying $*_k$ componentwise. Using the \inserted[typed]{composite} equivalence
\[ (119)\qquad F\underset{\deleted{\xout{\illegible{1}}}}{\hand{.}}\hand{\mapsto}(F.)^o :\ (D^-(A^\wedge))^o \hand{\approxeq} D^+((A^\wedge)^o) \hand{\approxeq} D^+(B^\wedge_{\hand{\mathcal{M}}}) \quad , \text{ with } \hand{\mathcal{M}} = \hand{\mathcal{N}}^o \ , \]
we get the \over[typed]{o}{t}autological duality isomorphism
\[ (120)\qquad (F.\overset{L}{*}_k L.')^o \hand{\simeq} RHom_k(L.',(F.)^o) \quad , \]
where the expression \deleted{\xout{\illegible{1}}} $\overset{\over{\unsure{P}}{L}}{\over{*}{*}}_k$ in the first member is relative to the pair
$(A,\hand{\mathcal{N}})$ , whereas the expression $RHom_k$ in the second is relative to the
dual pair $(B,\hand{\mathcal{M}})$. Symmetrically, we get
\[ (120')\qquad (RHom_k(L.,F^{\cdot}))^o \hand{\simeq} L.\overset{L}{*}_k (F.)^o \quad , \]
which is essentially the inverse isomorphism of (120), but for the pair
$(B,\hand{\mathcal{M}})$ instead of $(A,\hand{\mathcal{N}})$.

We still should duali\over[typed]{z}{s}e the functor \inserted[typed]{$R\over{I}{\Gamma}_A$} (116)(defined by (114)), which
we do recalling that $\over{I}{\Gamma}_A$ is just the \inserted[typed]{inverse limit} functor $\underset{A^o}{\underset{\hand{\longleftarrow}}{\lim}}$ , which is dual to \xout{th}
the direct limit functor $\underset{A^o}{\underset{\hand{\longrightarrow}}{\lim}}$, thus ``integ\over[typed]{g}{r}ation'' of \hand{$\mathcal{N}$}-valued
presheaves is just \inserted{(at least morally)} the total left derived functor of the latter, and
may be denoted by
\[ (121)\qquad L\,\underset{A^o}{\underset{\hand{\longrightarrow}}{\lim}} \quad , \]
while using for $R\over{I}{\Gamma}_A$ the equivalent notation, dual to (121)
\[ (12\over{2}{1}\hand{'})\qquad R\,\underset{A^o}{\underset{\hand{\longleftarrow}}{\lim}} = R\over{I}{\Gamma}_A \ . \]
I am not wholly happy though with the purely algebraic flavor of these
notations, not really suggestive of the manifold geometric intuition

\margin{113}
